\section{Numerical Experiment Setup}

\subsection{Test Problem}
\label{subsec:test_problem}

To verify and compare the two PINN formulations presented in Section~2, we construct a test problem based on the \emph{method of manufactured solutions}, which allows the error to be computed against a known exact solution.

We consider the case $d=1$, $\Omega = (0,1)$, $T=1$, and set the diffusion coefficient to $\nu = 0.1$. For the initial verification, the nonlinear term is chosen as $f(y) = y^3$.

We prescribe a reference state and control (the manufactured exact solution) as follows:
\begin{equation}
\label{eq:manufactured_y}
y^*(x,t) = \sin(\pi x)\, e^{-t},
\end{equation}
\begin{equation}
\label{eq:manufactured_u}
u^*(x,t) = y^*(x,t)\big|_{x\in\partial\Omega} = 0, \quad \forall t \in (0,T),
\end{equation}
since $\sin(\pi x) = 0$ at $x=0$ and $x=1$. The corresponding initial condition follows directly:
\begin{equation}
y_0(x) = y^*(x,0) = \sin(\pi x).
\end{equation}

Since the state equation \eqref{eq:state} contains no source term, the function $y^*$ above does not, in general, exactly satisfy the equation with $f(y)=y^3$. We therefore add an artificial source term $g(x,t)$ to the right-hand side of the state equation, obtained by substituting $y^*$ directly into the equation:
\begin{equation}
\label{eq:manufactured_source}
g(x,t) := \partial_t y^* - \nu \Delta y^* + f(y^*).
\end{equation}
With this source term, $y^*$ is the exact solution of the augmented state equation. The desired trajectory $y_d$ in the cost functional \eqref{eq:cost} is chosen to coincide with $y^*$, i.e., $y_d \equiv y^*$, so that $(y^*, u^*)$ is exactly the optimal solution of the control problem, with optimal cost value $J^* = 0$.

\subsection{Neural Network Configuration}
\label{subsec:network_setup}

Both formulations (direct and indirect) use fully-connected feedforward neural networks with $\tanh$ activation, taking $(x,t)$ as input and producing the corresponding scalar output.

\begin{table}[h!]
\centering
\caption{Neural network architectures used in the experiments.}
\label{tab:architecture}
\begin{tabular}{lccc}
\hline
Network & Hidden layers & Neurons/layer & Activation \\
\hline
$Y_\theta$ (state) & 4 & 50 & $\tanh$ \\
$U_\psi$ (control) & 3 & 30 & $\tanh$ \\
$\Lambda_\phi$ (adjoint, indirect method only) & 4 & 50 & $\tanh$ \\
\hline
\end{tabular}
\end{table}

For the indirect formulation, the distance function $d(x)$ used in the hard-constraint construction \eqref{eq:hard_constraint} is chosen as
\begin{equation}
d(x) = x(1-x),
\end{equation}
which satisfies $d(x) = 0$ for $x \in \{0,1\} = \partial\Omega$ and $d(x) > 0$ for $x \in (0,1)$.

Network weights are initialized using Xavier (Glorot uniform) initialization.

\subsection{Collocation Points and Sampling Strategy}

Training points are sampled using Latin Hypercube Sampling (LHS) over the corresponding domains:
\begin{itemize}
    \item $N_{\text{int}} = 10\,000$ interior points in $Q = (0,1)\times(0,1)$;
    \item $N_{\text{bc}} = 2\,000$ boundary points on $\Sigma = \{0,1\}\times(0,1)$;
    \item $N_{\text{ic}} = 500$ points for the initial condition at $t=0$;
    \item $N_T = 500$ points for the terminal condition at $t=T$ (used only in the indirect method, for the condition $\lambda(x,T)=0$).
\end{itemize}
These points are resampled every $500$ epochs to reduce overfitting to a fixed set of collocation points.

\subsection{Training Parameters}

Training is carried out in two stages:
\begin{enumerate}
    \item \textbf{Stage 1}: optimization with the Adam algorithm, initial learning rate $10^{-3}$, decayed exponentially (decay factor $0.9$ every $1\,000$ epochs), run for $20\,000$ epochs.
    \item \textbf{Stage 2}: fine-tuning with the L-BFGS algorithm until convergence or a maximum of $5\,000$ iterations, to improve accuracy once Adam has brought the solution near a minimum.
\end{enumerate}

Table~\ref{tab:hyperparams} summarizes the loss weights used for each method.

\begin{table}[h!]
\centering
\caption{Loss term weights for each formulation.}
\label{tab:hyperparams}
\begin{tabular}{lcc}
\hline
Term & Direct PINN & Indirect PINN \\
\hline
$w_{\text{res}}$ (state PDE residual) & $1$ & $1$ \\
$w_{\text{bc}}$ (state boundary condition, soft) & $10$ & --- (hard-constrained) \\
$w_{\text{ic}}$ (initial condition) & $10$ & $10$ \\
$w_\lambda$ (adjoint PDE residual) & --- & $1$ \\
$w_{\text{st}}$ (stationarity condition) & --- & $1$ \\
$w_T$ (adjoint terminal condition) & --- & $10$ \\
\hline
\end{tabular}
\end{table}

The Tikhonov regularization parameter is fixed at $\alpha = 10^{-2}$ throughout all experiments. Each configuration is trained over $5$ independent runs with different random seeds to assess the stability of the results.
